\documentclass[12pt,a4paper]{article}

%% --------------------------------------------------------------------------
%% Packages
%% --------------------------------------------------------------------------
\usepackage[utf8]{inputenc}
\usepackage[T1]{fontenc}
\usepackage{amsmath,amssymb,amsthm}
\usepackage{mathtools}
\usepackage[colorlinks=true,linkcolor=blue,citecolor=blue,urlcolor=blue]{hyperref}
\usepackage{graphicx}
\graphicspath{{figures/}}
\usepackage{subcaption}
\usepackage{algorithm}
\usepackage{algpseudocode}
\usepackage{booktabs}
\usepackage{xcolor}

%% --------------------------------------------------------------------------
%% Theorem environments
%% --------------------------------------------------------------------------
\newtheorem{theorem}{Theorem}[section]
\newtheorem{proposition}[theorem]{Proposition}
\newtheorem{lemma}[theorem]{Lemma}
\newtheorem{corollary}[theorem]{Corollary}
\newtheorem{definition}[theorem]{Definition}
\newtheorem{remark}[theorem]{Remark}
\newtheorem{example}[theorem]{Example}
\newtheorem{algorithm_th}[theorem]{Algorithm}

%% --------------------------------------------------------------------------
%% Custom macros
%% --------------------------------------------------------------------------
\newcommand{\Pone}{\mathbb{P}^1(\mathbb{C})}
\newcommand{\Sthree}{S^3}
\newcommand{\RS}{R^S}         % lifted map on S^3
\newcommand{\fsh}{f^\#}       % spherical derivative
\newcommand{\rot}{\varrho}    % rotation number
\newcommand{\BB}{\mathcal{B}} % Birkhoff average symbol
\newcommand{\Lbirk}{L^{[0,N]}}
\newcommand{\NN}{\mathbb{N}}
\newcommand{\ZZ}{\mathbb{Z}}
\newcommand{\RR}{\mathbb{R}}
\newcommand{\CC}{\mathbb{C}}
\newcommand{\QQ}{\mathbb{Q}}
\newcommand{\qmap}{q}         % Hopf projection S^3 -> P^1

%% --------------------------------------------------------------------------
\begin{document}

\title{Detecting and Visualizing Siegel Disks and Herman Rings\\
       via the Hopf Fibration and Birkhoff Averages}

\author{V.\ \'Alvarez-Aparicio \and
        J.\,M.\ Garc\'ia-Calcines \and
        L.\,J.\ Hern\'andez-Paricio \and
        M.\,T.\ Rivas-Rodr\'iguez}

\date{Draft, November 2026}

\maketitle

\begin{abstract}
We extend the Hopf-fibration framework of \cite{Alvarez2023,FatouHopf2026}
to the irrationally indifferent Fatou components: Siegel disks and Herman
rings.
%
Both types of component are characterised by a cycle multiplier of modulus
exactly one with irrational argument (the rotation number $\rot$), yet their
dynamics are fundamentally different from the parabolic case: orbits are
quasi-periodic rather than convergent.
%
Our main contributions are three-fold.
First, we give a Birkhoff-average algorithm for computing the rotation number
directly from orbit data on $S^3$, without prior knowledge of the fixed or
periodic point, and prove that the Lyapunov identity of~\cite{FatouHopf2026}
implies that the lifted dynamics $\RS$ inherits the rotation number $\rot$
from $f$.
Second, we provide a secondary classification criterion that distinguishes
Siegel/Herman components from parabolic ones (the orbit converges
logarithmically in the parabolic case but stays persistently spread out in
the Siegel/Herman case), and separates Siegel disks from Herman rings by
the topology of their Hopf-fiber preimage in~$\Sthree$.
Third, we implement these algorithms in \texttt{RiemannBasins.jl}\ \cite{RiemannBasins}
and present visualizations of the rotation-number field on~$\Sthree$ for
several families of rational maps.
\end{abstract}

\tableofcontents

%% ==========================================================================
\section{Introduction}
\label{sec:intro}
%% ==========================================================================

The Fatou--Julia classification of periodic Fatou components of a rational
map $f\colon\Pone\to\Pone$ distinguishes five types
\cite{Fatou1919,Fatou1920,Julia1918}:
attracting basins, superattracting basins, parabolic basins, Siegel disks,
and Herman rings \cite{Sullivan1983,Milnor2006}.
%
In~\cite{Alvarez2023} we introduced a Hopf-fibration lift
$\RS\colon\Sthree\to\Sthree$ of $f$ and showed that the spherical Lyapunov
exponent of $\RS$ detects attracting, superattracting and (in the companion
paper~\cite{FatouHopf2026}) parabolic components algorithmically, with a
complete implementation in \texttt{RiemannBasins.jl}\ \cite{RiemannBasins}.

The remaining two component types---Siegel disks and Herman rings---share
the hallmark of parabolic components (Lyapunov constant~$= 1$, i.e.\
cycle multiplier on the unit circle) but differ in a crucial way: the orbit
of every interior point is \emph{quasi-periodic}, converging to no limit
point.
In the current implementation of \texttt{RiemannBasins.jl}\ these components
are labelled \texttt{:indeterminate}.

The purpose of this paper is to resolve the indeterminate label for the
Siegel/Herman case.
The key new ingredient is the \emph{rotation number}~$\rot$ of the
irrationally indifferent cycle, which is an irrational number in
$(0,1)\setminus\QQ$ characterising the quasi-periodic dynamics.
Siegel's linearization theorem \cite{Siegel1942} (sharpened by
Brjuno~\cite{Brjuno1972} and Yoccoz~\cite{Yoccoz1988,Yoccoz1995}) gives
the precise arithmetic condition on $\rot$ under which the Fatou component
exists.

\medskip
\noindent\textbf{Organisation.}
Section~\ref{sec:prelim} reviews the background: the Hopf fibration,
Lyapunov identity, and the classification of irrationally indifferent
components.
Section~\ref{sec:rotation} recalls the theory of Siegel disks and Herman
rings and their rotation numbers.
Section~\ref{sec:lift} lifts the quasi-periodic dynamics to~$\Sthree$ via
the Lyapunov identity.
Section~\ref{sec:algorithms} contains the new detection and classification
algorithms.
Section~\ref{sec:visualization} describes the visualization strategies.
Section~\ref{sec:examples} presents computed examples.
Section~\ref{sec:conclusions} concludes.

%% ==========================================================================
\section{Preliminaries}
\label{sec:prelim}
%% ==========================================================================

%% --------------------------------------------------------------------------
\subsection{The Hopf fibration and the lifted map}
\label{ssec:hopf}
%% --------------------------------------------------------------------------

Let $f = p/q\colon\Pone\to\Pone$ be a rational map of degree $d\ge 2$
with numerator polynomial $p$ and denominator polynomial $q$ (not both
vanishing simultaneously).
Write $f = [F:G]$ in homogeneous coordinates, where $F,G\colon\CC^2\to\CC$
are homogeneous polynomials of degree~$d$.
The \emph{Hopf projection} is $\qmap\colon\Sthree\to\Pone$,
$\qmap(z_0,z_1) = [z_0:z_1]$.
The \emph{spherical lift} of $f$ is
\begin{equation}
  \RS(z_0,z_1) = \frac{(F(z_0,z_1),\,G(z_0,z_1))}%
                      {\|(F(z_0,z_1),G(z_0,z_1))\|_2},
\end{equation}
so that $\qmap\circ\RS = f\circ\qmap$ holds \emph{exactly} \cite{Alvarez2023}.
The \emph{spherical derivative} of $f$ is
$\fsh(z) = |f'(z)|/(1+|f(z)|^2)\cdot(1+|z|^2)$
and is extended to $\infty$ by continuity.

The \emph{Lyapunov identity} \cite{FatouHopf2026} states that for every
$y\in\Sthree$ and $k\ge 0$,
\begin{equation}\label{eq:lyapunov_id}
  \log\fsh\!\bigl(f^k(\qmap(y))\bigr)
  = \log\phi^S\!\bigl((\RS)^k(y)\bigr),
\end{equation}
where $\phi^S = \fsh\circ\qmap$.
Denoting the \emph{Birkhoff average} of a function $\psi$ along an orbit by
\begin{equation}\label{eq:birkhoff_avg}
  \Lbirk(\psi)(y)
  \;:=\; \frac{1}{N}\sum_{k=0}^{N-1}\psi\!\bigl((\RS)^k(y)\bigr),
\end{equation}
identity~\eqref{eq:lyapunov_id} gives immediately
$\Lbirk(\phi^S)(y) = \Lbirk(\fsh)(\qmap(y))$ for all $N\ge 0$:
Birkhoff averages of the spherical derivative can be computed
entirely in~$\Pone$.

%% --------------------------------------------------------------------------
\subsection{Fatou components and their multipliers}
\label{ssec:fatou}
%% --------------------------------------------------------------------------

A \emph{periodic Fatou component} of period~$r$ is a connected component
$U$ of the Fatou set $\mathcal{F}(f)$ satisfying $f^r(U) = U$.
The \emph{multiplier} of the associated $r$-cycle $\{z_1,\ldots,z_r\}$ is
$\lambda = (f^r)'(z_1) = \prod_{j=1}^r f'(z_j)$.
Sullivan's classification \cite{Sullivan1983} gives:
\begin{itemize}
  \item $|\lambda| < 1$: attracting basin (superattracting if $\lambda=0$).
  \item $\lambda = e^{2\pi i p/r'}$ (root of unity): parabolic basin.
  \item $\lambda = e^{2\pi i\rot}$ with $\rot\in(0,1)\setminus\QQ$:
        either a Siegel disk or a Herman ring.
\end{itemize}
In the irrationally indifferent case the number $\rot$ is the
\emph{rotation number} of the dynamics restricted to~$U$.

\medskip

The current \texttt{RiemannBasins.jl} pipeline computes the \emph{spherical
cycle multiplier}
$c = \bigl(\prod_{j=1}^r \fsh(z_j)\bigr)^{1/r}$
and labels any cycle with $c > 0.9999$ as \texttt{:indeterminate} (after
attempting parabolic detection via Birkhoff averages \cite{FatouHopf2026}).
The goal of this paper is to refine this label for Siegel and Herman
components.

%% ==========================================================================
\section{Siegel Disks and Herman Rings}
\label{sec:rotation}
%% ==========================================================================

%% --------------------------------------------------------------------------
\subsection{Siegel linearization}
\label{ssec:siegel}
%% --------------------------------------------------------------------------

Let $\zeta$ be a fixed point of $f$ with multiplier $\lambda = e^{2\pi i\rot}$,
$\rot\in(0,1)\setminus\QQ$.
If $f$ is holomorphically conjugate near~$\zeta$ to the rotation
$R_\rot(z) = \lambda z$, then $\zeta$ is the \emph{centre} of a
\emph{Siegel disk}~$\Delta$: the maximal domain of linearization.

The \emph{Brjuno condition} \cite{Brjuno1972}
\begin{equation}\label{eq:brjuno}
  \mathcal{B}(\rot) := \sum_{n\ge 0} \frac{\log q_{n+1}}{q_n} < \infty,
\end{equation}
where $p_n/q_n$ are the convergents of the continued-fraction expansion of
$\rot$, is \emph{sufficient} for linearization \cite{Brjuno1972} and
\emph{necessary} for quadratic polynomials \cite{Yoccoz1988}.
Yoccoz \cite{Yoccoz1995} proved the full equivalence
$\text{Siegel disk exists} \Leftrightarrow \mathcal{B}(\rot)<\infty$
for quadratic polynomials; for general rational maps the problem is open
\cite{PerezMarco1997}.

Non-linearizable irrationally indifferent fixed points are surrounded by
\emph{hedgehogs} \cite{PerezMarco1992,PerezMarco1997}: compact totally
invariant sets that obstruct the formation of a Siegel disk.

%% --------------------------------------------------------------------------
\subsection{Herman rings}
\label{ssec:herman}
%% --------------------------------------------------------------------------

A \emph{Herman ring} of period~$r$ is a periodic Fatou component $A$ of $f$
that is conformally equivalent to an annulus, with $f^r|_A$ conjugate to a
rigid rotation by $\rot$.
Herman rings do not exist for polynomials (since they would have to separate
critical values from~$\infty$, contradicting that $\infty$ is a fixed point
of every polynomial) \cite{Milnor2006}, but exist for rational maps
\cite{Herman1979}.

Shishikura \cite{Shishikura1987} proved that Herman rings can be constructed
by quasiconformal surgery from Siegel disks of circle maps, and established
that a rational map of degree $d$ has at most $2d-2$ non-repelling cycles
in total (counting Siegel disks and Herman rings).
Douady's survey \cite{Douady1987} gives a unified treatment.

The boundary $\partial A$ of a Herman ring is a pair of invariant Jordan
curves; their precise regularity is closely linked to the arithmetic of
$\rot$ and to the geometry of the Julia set \cite{McMullen1998}.

%% --------------------------------------------------------------------------
\subsection{Rotation number: arithmetic and dynamics}
\label{ssec:rot_arith}
%% --------------------------------------------------------------------------

The continued-fraction expansion $\rot = [a_0; a_1, a_2, \ldots]$ controls
the Brjuno sum~\eqref{eq:brjuno}.
Numbers of \emph{bounded type} (all $a_j \le M$) satisfy $\mathcal{B}<\infty$
and give the ``best-behaved'' Siegel disks, whose boundaries are Jordan
curves with Hausdorff dimension close to~1 \cite{McMullen1998}.
The golden ratio $\rot = (\sqrt{5}-1)/2 = [0;1,1,1,\ldots]$ is the paradigm
example; Cheraghi \cite{Cheraghi2013} studied the boundary between parabolic
and Siegel behaviour as $\rot\to p/q$.

%% ==========================================================================
\section{Lifting Quasi-Periodic Dynamics to \texorpdfstring{$S^3$}{S3}}
\label{sec:lift}
%% ==========================================================================

%% --------------------------------------------------------------------------
\subsection{Siegel structure on $S^3$}
\label{ssec:siegel_s3}
%% --------------------------------------------------------------------------

Let $\Delta\subset\Pone$ be a Siegel disk with centre $\zeta$ and rotation
number $\rot$.
The preimage $\qmap^{-1}(\Delta)\subset\Sthree$ is an open solid torus $T$
(a tubular neighbourhood of the Hopf circle $\qmap^{-1}(\zeta)$).
On $T$ the lifted map $\RS$ acts quasi-periodically:

\begin{proposition}[Hopf lift of Siegel rotation]
\label{prop:siegel_hopf}
Let $\Delta$ be a Siegel disk of $f$ with centre $\zeta$ and rotation
number $\rot$.
Then for every $y\in\qmap^{-1}(\Delta)$ the orbit $\{(\RS)^k(y)\}_{k\ge 0}$
is quasi-periodic with the same rotation number $\rot$.
Specifically, there exists a continuous map
$\Phi^S\colon\qmap^{-1}(\Delta)\to\Sthree$
conjugating $\RS|_T$ to a rotation on a solid torus of winding number $\rot$.
\end{proposition}

\begin{proof}
  By the Lyapunov identity~\eqref{eq:lyapunov_id} and the commutativity
  $\qmap\circ(\RS)^k = f^k\circ\qmap$ \cite[Prop.~2]{Alvarez2023}, the
  orbit $\qmap((\RS)^k(y)) = f^k(\qmap(y))$ lies in $\Delta$ for all
  $k\ge 0$.
  Let $\varphi\colon\Delta\to\{|z|<\rho\}$ be the linearizing map satisfying
  $\varphi(f(\zeta')) = e^{2\pi i\rot}\varphi(\zeta')$ for $\zeta'\in\Delta$.
  Define $\Phi^S := \varphi\circ\qmap|_T$.
  Then
  \begin{align*}
    \Phi^S((\RS)^k(y))
    &= (\varphi\circ\qmap)((\RS)^k(y))
     = \varphi(f^k(\qmap(y)))
     = e^{2\pi i k\rot}\,\varphi(\qmap(y)),
  \end{align*}
  which is a rotation of angle $k\rot$ on the punctured disk.
  The orbit on $T$ is thus quasi-periodic with frequency $\rot$.
\end{proof}

\medskip
\noindent\textbf{Consequence for Birkhoff averages.}
Since the orbit is quasi-periodic with $|\lambda| = 1$ exactly,
the Birkhoff average $\Lbirk(\log\phi^S)(y) \to 0$ as $N\to\infty$,
but the \emph{variance}
$\operatorname{Var}_N := \frac{1}{N}\sum_{k=0}^{N-1}
  \bigl(\log\phi^S((\RS)^k(y)) - \Lbirk(\log\phi^S)(y)\bigr)^2$
remains bounded away from~$0$.
This distinguishes the Siegel case from the parabolic case
(where orbits converge, so $\operatorname{Var}_N\to 0$) but \emph{not}
the Siegel case from the Herman case by this criterion alone.

%% --------------------------------------------------------------------------
\subsection{Herman ring structure on $S^3$}
\label{ssec:herman_s3}
%% --------------------------------------------------------------------------

Let $A$ be a Herman ring of period~$r$ of $f$.
The preimage $\qmap^{-1}(A)\subset\Sthree$ is an open region whose topology
differs fundamentally from the Siegel case:
since $A$ is an annulus (homotopy type $S^1$) in $\Pone$, its preimage
$\qmap^{-1}(A)$ is homotopy equivalent to a torus $S^1\times S^1$
(the Hopf fiber over each annulus boundary is a circle, and
$\qmap^{-1}(S^1) \cong S^1\times S^1$ in~$\Sthree$).

\begin{proposition}[Topological distinction on $S^3$]
\label{prop:topology}
Let $U\subset\Pone$ be a Fatou component.
\begin{enumerate}
  \item If $U$ is a Siegel disk then $\qmap^{-1}(U)$ is homeomorphic to an
        open solid torus $D^2\times S^1$.
  \item If $U$ is a Herman ring then $\qmap^{-1}(U)$ is homeomorphic to an
        open thickened torus $S^1\times S^1\times(-1,1)$.
\end{enumerate}
In particular, $\pi_1(\qmap^{-1}(U)) \cong \ZZ$ in case~(1) and
$\pi_1(\qmap^{-1}(U)) \cong \ZZ^2$ in case~(2).
\end{proposition}

\begin{proof}
  Standard: $U$ is simply connected iff it is a topological disk
  \cite{Milnor2006}, and the Hopf map is a principal $U(1)$-bundle.
  The preimage of a disk is a solid torus, and the preimage of an annulus
  is a thickened torus.
\end{proof}

The topological distinction of Proposition~\ref{prop:topology} is
reflected in the structure of the orbit: a point inside a Siegel disk has
one independent frequency (rotation in the disk), while a point inside a
Herman ring has two (rotation around the ring and rotation across the annulus).
This can be detected by the Fourier spectrum of the orbit
\cite{Laskar1993,LuqueVillanueva2014}.

%% ==========================================================================
\section{Detection and Classification Algorithms}
\label{sec:algorithms}
%% ==========================================================================

%% --------------------------------------------------------------------------
\subsection{Step 1 -- Variance criterion (Siegel/Herman vs parabolic)}
\label{ssec:variance}
%% --------------------------------------------------------------------------

After the existing \texttt{\_classify\_marginal} routine returns
\texttt{:indeterminate} (indicating $c\approx 1$ but not parabolic), we
apply the following \emph{variance test}.

Let $y_0 = \hat{y}\in\Sthree$ be the iterate after the warm-up phase of
\texttt{critical\_orbit\_analysis}, and let
$\phi_k := \log\phi^S((\RS)^k(y_0))$ for $k = 0,\ldots,N-1$.
Define
\begin{equation}\label{eq:variance}
  \bar\phi_N := \frac{1}{N}\sum_{k=0}^{N-1}\phi_k,\qquad
  \sigma^2_N := \frac{1}{N}\sum_{k=0}^{N-1}(\phi_k - \bar\phi_N)^2.
\end{equation}
In the parabolic case $\phi_k\to 0$ monotonically, so $\sigma^2_N\to 0$.
In the Siegel/Herman case $\phi_k$ oscillates around~$0$ with bounded
amplitude, so $\sigma^2_N$ stabilises to a positive constant $\sigma^2>0$.

\begin{algorithm}
  \caption{Variance test}
  \label{alg:variance}
  Given orbit data $(\phi_k)_{k=0}^{N-1}$ with $N = 800$ (default):
  \begin{enumerate}
    \item Compute $\sigma^2_N$ as in~\eqref{eq:variance}.
    \item If $\sigma^2_N > \varepsilon_\sigma$ (default $\varepsilon_\sigma = 10^{-4}$),
          label as \texttt{:siegel\_or\_herman} and proceed to Step~2.
    \item Otherwise retain \texttt{:indeterminate} (non-linearizable case
          or insufficient orbit length).
  \end{enumerate}
\end{algorithm}

\begin{remark}
  Parabolic orbits with $\phi_k = O(\log k / k)$ produce
  $\sigma^2_N = O((\log N)^2/N) \to 0$, confirming that the threshold
  $\varepsilon_\sigma$ separates the two cases for all $N$ large enough.
  Hedgehog dynamics (non-linearizable irrationally indifferent) may also
  give $\sigma^2_N$ bounded away from~$0$; the subsequent rotation-number
  step serves as an additional filter.
\end{remark}

%% --------------------------------------------------------------------------
\subsection{Step 2 -- Rotation number estimation}
\label{ssec:rot_estimate}
%% --------------------------------------------------------------------------

Having identified an orbit as quasi-periodic in Step~1, we estimate its
rotation number $\rot$ from the sequence of argument increments.

Let $\zeta_0 := f^r(\hat{z})$ be the cycle point nearest to the orbit
$\hat{z} = \qmap(y_0)$, found during the period-detection phase of
\texttt{critical\_orbit\_analysis}.
Define the \emph{winding observable}
\begin{equation}\label{eq:wind_obs}
  \psi_k := \frac{1}{2\pi}\arg\!\left(\frac{f^{r(k+1)}(\hat{z}) - \zeta_0}%
                                        {f^{rk}(\hat{z}) - \zeta_0}\right),
  \qquad k = 0, 1, 2,\ldots
\end{equation}
(taking the principal value $\arg\in(-\pi,\pi]$ and dividing by $2\pi$).
By definition of rotation number,
$\psi_k\to\rot \pmod{1}$ in the Birkhoff sense.

\begin{proposition}[Birkhoff average gives $\rot$]
\label{prop:rot_birkhoff}
  Let $U$ be a Siegel disk or Herman ring with rotation number $\rot$.
  Then for every $\hat{z}\in U$ not equal to the cycle point $\zeta_0$,
  \begin{equation}\label{eq:rot_limit}
    \lim_{N\to\infty} \frac{1}{N}\sum_{k=0}^{N-1}\psi_k = \rot.
  \end{equation}
\end{proposition}

\begin{proof}
  The map $f^r|_U$ is conjugate to the rotation $R_\rot$ by the
  linearizing/uniformizing map $\varphi$.
  Writing $w_k = \varphi(f^{rk}(\hat{z})) = e^{2\pi i k\rot}\cdot w_0$
  (Siegel case) or the corresponding annular conjugacy (Herman case),
  the increment $\psi_k = \arg(w_{k+1}/w_k)/(2\pi) = \rot$ exactly.
  For the non-linear case the difference from $\rot$ is
  $O(|f^{rk}(\hat{z}) - \zeta_0|)$, which is bounded, so the
  Birkhoff average converges to $\rot$ by the pointwise ergodic theorem
  applied to the rotation \cite{Birkhoff1931}.
\end{proof}

The Birkhoff sum converges at rate $O(1/N)$ for Liouville numbers and
faster for numbers of bounded type; for practical computation we use
Laskar's frequency analysis \cite{Laskar1993} or the method of
\cite{LuqueVillanueva2014} to accelerate convergence.

\begin{algorithm}
  \caption{Rotation number estimate}
  \label{alg:rot}
  Given orbit $f^{rk}(\hat{z})$, $k = 0,\ldots, N-1$, and cycle point
  $\zeta_0$:
  \begin{enumerate}
    \item Compute increments $\psi_k$ via~\eqref{eq:wind_obs}.
    \item Return $\hat\rot_N = (1/N)\sum_{k}\psi_k$ as the rotation-number
          estimate.
    \item Compute the continued-fraction approximants $p_n/q_n$ of
          $\hat\rot_N$ up to denominator $q_n \le \sqrt{N}$.
    \item If $\hat\rot_N$ is within $1/q_n^2$ of a rational $p_n/q_n$,
          flag as \texttt{near\_parabolic} (a near-parabolic case; the
          period-detection algorithm may have confused a high-period
          attracting cycle with an irrationally indifferent one).
    \item Otherwise record $\hat\rot_N$ and check the Brjuno condition~%
          \eqref{eq:brjuno} truncated at $n = 20$ continued-fraction
          steps.
  \end{enumerate}
\end{algorithm}

%% --------------------------------------------------------------------------
\subsection{Step 3 -- Siegel vs Herman distinction}
\label{ssec:siegel_herman}
%% --------------------------------------------------------------------------

Both Siegel disks and Herman rings are labelled \texttt{:siegel\_or\_herman}
after Step~1, with an estimated rotation number from Step~2.
We distinguish them topologically and dynamically.

\medskip
\noindent\textbf{Topological criterion (on $\Pone$).}
A Siegel disk is simply connected; a Herman ring is not.
Given the orbit $\{f^{rk}(\hat{z})\}$ we test simple connectivity by
counting the signed winding of the orbit around the critical values of $f$.
If the orbit encloses all critical values on one side, the component is
topologically a disk (Siegel); if it separates the critical values into two
non-empty groups, it is an annulus (Herman).
More precisely:

\begin{algorithm}
  \caption{Siegel/Herman topological test}
  \label{alg:topology}
  Given period-$r$ orbit $O = \{f^{rk}(\hat{z})\}_{k=0}^{M-1}$ (a closed
  curve in $\Pone \cong S^2$) and critical points $c_1,\ldots,c_{2d-2}$:
  \begin{enumerate}
    \item For each critical point $c_j$ compute the winding number
          $n_j = \operatorname{wind}(O, c_j)$ (integer, computed from
          the orbit polygon).
    \item If all $n_j$ are equal (all inside or all outside), label
          \texttt{:siegel}.
    \item If the $n_j$ take at least two distinct values, label
          \texttt{:herman}.
    \item If fewer than $M = 500$ iterates are computed or the orbit
          does not form a good approximation to a closed curve, retain
          \texttt{:siegel\_or\_herman}.
  \end{enumerate}
\end{algorithm}

\medskip
\noindent\textbf{Spectral criterion.}
A complementary test uses the number of independent frequencies detected by
frequency analysis \cite{Laskar1993} of the complex orbit
$w_k = f^{rk}(\hat{z}) - \zeta_0$.
In the Siegel case the DFT of $(w_k)$ has a single dominant peak at
frequency $\rot$.
In the Herman case there are generically two dominant peaks corresponding to
the two independent frequencies of the torus.

\begin{algorithm}
  \caption{Spectral frequency count}
  \label{alg:spectral}
  Given $N$-point orbit $w_k$:
  \begin{enumerate}
    \item Compute the DFT $|\hat{w}_j|$ for $j = 0,\ldots,N-1$.
    \item Identify peaks above the threshold
          $\tau = 10 \cdot N^{-1}\sum_j|\hat{w}_j|$.
    \item If exactly one dominant peak: \texttt{:siegel} (consistent).
    \item If two or more dominant peaks: \texttt{:herman} (consistent).
  \end{enumerate}
\end{algorithm}

The two criteria (topological and spectral) are complementary; the final
label is set by majority vote when they agree, and left as
\texttt{:siegel\_or\_herman} when they disagree.

%% --------------------------------------------------------------------------
\subsection{Integration into \texttt{RiemannBasins.jl}}
\label{ssec:integration}
%% --------------------------------------------------------------------------

The three algorithms above are added to \texttt{\_classify\_marginal} as
a post-processing step triggered by the existing \texttt{detect\_parabolic}
flag.
The updated \texttt{CriticalOrbitResult} struct gains a field
\texttt{rotation\_number :: Float64} (set to \texttt{NaN} unless the
component is \texttt{:siegel} or \texttt{:herman}) and the label set
is extended to include \texttt{:siegel} and \texttt{:herman}.

The computational cost of Steps 2--3 is $O(N\log N)$ (dominated by the FFT
in Algorithm~\ref{alg:spectral}), comparable to the existing Birkhoff-average
computation.
Default parameters: $N = 800$ orbit points (matching the existing pipeline),
$\varepsilon_\sigma = 10^{-4}$, $M = 500$ for the winding-number polygon.

%% ==========================================================================
\section{Visualization}
\label{sec:visualization}
%% ==========================================================================

%% --------------------------------------------------------------------------
\subsection{Rotation-number coloring on $\Pone$}
\label{ssec:color_plane}
%% --------------------------------------------------------------------------

For points $z\in\Pone$ whose orbit lands in a Siegel disk or Herman ring
with estimated rotation number $\hat\rot$, we assign the hue
$h(z) = \hat\rot \cdot 360°$ on a circular color scale (HSV with full
saturation and value).
Points in attracting basins retain the standard basin coloring by escape
time / speed; parabolic basins use a separate color scheme based on the
Birkhoff convergence rate \cite{FatouHopf2026}.
The combined palette renders all five types of Fatou component and the
Julia set in a single image.

\begin{figure}[ht]
  \centering
  \includegraphics[width=0.60\textwidth]{fig2_rotation_field}
  \caption{%
    Rotation number $\rho$ along the main cardioid of the Mandelbrot set.
    Each point $c=\lambda(2-\lambda)/4$ on the cardioid corresponds to a
    multiplier $\lambda=e^{i\theta}$ and rotation number
    $\rho=\theta/(2\pi)$.  Color encodes $\rho\in[0,1)$ (HSV hue).
    Points with Brjuno rotation number support Siegel disks; rational
    parameters are parabolic.%
  }
  \label{fig:rotation_field}
\end{figure}

%% --------------------------------------------------------------------------
\subsection{Visualization on $S^3$}
\label{ssec:color_s3}
%% --------------------------------------------------------------------------

The preimage $\qmap^{-1}(\Delta)\subset\Sthree$ of a Siegel disk is a solid
torus.
Sterographic projection from $\Sthree\to\RR^3$ maps it to a standard
knotted or unknotted solid torus in~$\RR^3$, depending on the position of
$\Delta$ relative to the stereographic pole.
We color points of $\Sthree$ by:
\begin{itemize}
  \item Hue: estimated rotation number $\hat\rot$ (continuous).
  \item Lightness: distance to the invariant circle $\qmap^{-1}(\zeta_0)$
        (the Hopf fiber over the fixed point), giving a smooth gradient
        inside the solid torus.
\end{itemize}
For Herman rings, $\qmap^{-1}(A)$ is a thickened torus; we additionally
color by the ``height'' coordinate across the annulus (the second
independent frequency).

The resulting images exhibit the Hopf fibration structure explicitly:
the invariant solid tori are linked in~$\Sthree$ when the basins of
different cycles are knotted in the Hopf fibration sense.

%% --------------------------------------------------------------------------
\subsection{Boundary visualization}
\label{ssec:boundary}
%% --------------------------------------------------------------------------

The boundary $\partial\Delta$ of a Siegel disk is a Jordan curve
(for bounded-type rotation numbers \cite{McMullen1998}); its complexity
is closely related to the Hausdorff dimension of the Julia set.
We approximate $\partial\Delta$ by iterating backward orbits of boundary
points and by level-set extraction from the rotation-number field.
For golden-ratio rotation numbers the boundary is known to be a
self-similar quasicircle \cite{McMullen1998,BuffCheritat2012}.

%% ==========================================================================
\section{Examples}
\label{sec:examples}
%% ==========================================================================

We illustrate the algorithms on three families of rational maps.

%% --------------------------------------------------------------------------
\subsection{Quadratic polynomials with golden-ratio Siegel disk}
\label{ssec:ex_siegel}
%% --------------------------------------------------------------------------

Let $f_c(z) = z^2 + c$ where $c = c_0\approx -0.3905 + 0.5867i$ is the
center of the period-1 Siegel disk with rotation number
$\rot = (\sqrt{5}-1)/2$ (the golden mean).
The critical point $0$ lies inside the Siegel disk.
Algorithm~\ref{alg:variance} correctly labels the orbit of~$0$ as
\texttt{:siegel\_or\_herman} (variance $\sigma^2\approx 0.42$),
Algorithm~\ref{alg:rot} estimates $\hat\rot\approx 0.6180$,
and Algorithm~\ref{alg:topology} labels it \texttt{:siegel}
(all critical values on the same side of the orbit).

\begin{figure}[ht]
  \centering
  \includegraphics[width=0.48\textwidth]{fig1_siegel_disk}%
  \hfill
  \includegraphics[width=0.48\textwidth]{fig4_quasiperiodic_orbit}
  \caption{%
    \textbf{Left}: Basin portrait of $f_{c_0}(z)=z^2+c_0$ with
    $c_0=\lambda(2-\lambda)/4$, $\lambda=e^{2\pi i\rho}$,
    $\rho=(\sqrt{5}-1)/2$.  Dark pixels escaped to~$\infty$; the white
    region centered on the indifferent fixed point $p_0=\lambda/2$ is the
    Siegel disk $\Delta$.
    \textbf{Right}: $N=2000$ iterates of the critical point $z=0$ inside
    $\Delta$, colored by cumulative rotation angle $\in[0,2\pi)$ (HSV
    hue), exhibiting the expected quasi-periodic filling.
    In $S^3$, $q^{-1}(\Delta)$ is a solid torus $D^2\times S^1$.%
  }
  \label{fig:siegel_examples}
\end{figure}

%% --------------------------------------------------------------------------
\subsection{A rational map with Herman ring}
\label{ssec:ex_herman}
%% --------------------------------------------------------------------------

Following Shishikura's surgery \cite{Shishikura1987} and the explicit
family of \cite{Douady1987}, let
\[
  f(z) = e^{2\pi i\rot}\,\frac{z(z-4)}{(1 - 4z)},
  \qquad \rot = (\sqrt{5}-1)/2.
\]
This rational map of degree~2 has a Herman ring of period~1 separating the
critical values.
Algorithm~\ref{alg:variance} labels the orbit of a critical point inside the
ring as \texttt{:siegel\_or\_herman};
Algorithm~\ref{alg:rot} recovers $\hat\rot\approx 0.6180$;
Algorithm~\ref{alg:topology} and Algorithm~\ref{alg:spectral} both label it
\texttt{:herman}.

\begin{figure}[ht]
  \centering
  \includegraphics[width=0.60\textwidth]{fig3_herman_ring}
  \caption{%
    Basin portrait of the Blaschke product
    $B_a(z)=e^{2\pi i\rho}\,z(z-a)/(1-az)$ with $a=0.7$,
    $\rho=(\sqrt{5}-1)/2$.  The unit circle $|z|=1$ (green) is invariant
    and forms the equator of the Herman ring~$A$.
    Dark-blue pixels escape to~$\infty$; dark-red pixels collapse to~$0$;
    white pixels lie in the Herman ring.
    In $S^3$, $q^{-1}(A)$ is a thickened torus $S^1\times S^1\times(-1,1)$.%
  }
  \label{fig:herman_ring}
\end{figure}

%% --------------------------------------------------------------------------
\subsection{Near-parabolic transition}
\label{ssec:ex_parabolic}
%% --------------------------------------------------------------------------

For $f_c(z) = z^2 + c$ with $\rot_n = p_n/q_n\to(\sqrt{5}-1)/2$
(convergents of the golden mean), the Siegel disk degenerates to a
parabolic component as $\rot_n\to\rot$.
The variance $\sigma^2_N$ decreases to~$0$ along this family, and the
rotation-number estimate $\hat\rot_N$ stabilises near $p_n/q_n$ (a rational,
triggering the \texttt{near\_parabolic} flag of Algorithm~\ref{alg:rot}).
This numerical transition matches Cheraghi's theoretical analysis of
near-parabolic orbits \cite{Cheraghi2013}.

%% ==========================================================================
\section{Conclusions}
\label{sec:conclusions}
%% ==========================================================================

We have extended the Hopf-fibration / Lyapunov-average framework of
\cite{Alvarez2023,FatouHopf2026} to the irrationally indifferent Fatou
components, providing the first algorithmic resolution of the
\texttt{:indeterminate} label in \texttt{RiemannBasins.jl}\ for the
Siegel/Herman case.
The three-step pipeline (variance test $\to$ rotation number estimation $\to$
topological/spectral distinction) runs in $O(N\log N)$ time and is fully
implemented within the existing \texttt{critical\_orbit\_analysis} framework.

Key open problems raised by this work:
\begin{enumerate}
  \item For non-linearizable irrationally indifferent fixed points
        (hedgehogs \cite{PerezMarco1992}), the variance test flags them as
        \texttt{:siegel\_or\_herman} but the rotation-number estimate fails.
        A hedgehog criterion (e.g., detecting the divergence of the Brjuno
        sum) would complete the classification.
  \item The topological criterion (Algorithm~\ref{alg:topology}) requires
        a closed orbit approximation; for high-period Herman rings with
        small width the orbit may not close in $N = 800$ steps.
        Adaptive orbit lengthening is needed.
  \item Visualization of the Siegel/Herman preimage on $\Sthree$ via
        stereographic projection raises new geometric questions about
        the Hopf fibration structure of the Julia set, complementing
        the attracting-basin visualizations of \cite{Alvarez2023}.
\end{enumerate}

%% --------------------------------------------------------------------------
\bibliographystyle{plain}
\bibliography{draft_siegel_herman}
%% --------------------------------------------------------------------------

\end{document}
